\section{The tensorizable Hilbert benchmark}
\label{sec:tpc42-hilbert-large-sieve}

There are two logically different ways in which a Hilbert space can enter
a large--sieve argument.  One may first prove a scalar arithmetic estimate
and then apply a bounded synthesis map to its output, or one may insert a
vector coefficient before the arithmetic estimate is proved.  The first
operation is automatic and dimension free.  The second is not.  This
section records the exact distinction at the TPC--41 progression scale.

\subsection{Dimension--free postcomposition}

\begin{lemma}[Hilbert tensorization of a scalar operator]
\label{lem:tpc42-hilbert-tensorization}
Let \(I,J\) be finite sets and let \(T=(t_{ji})_{j\in J,i\in I}\) satisfy
\begin{equation}
 \sum_{j\in J}\left|\sum_{i\in I}t_{ji}z_i\right|^2
 \le C^2\sum_{i\in I}|z_i|^2
 \qquad(z\in\C^I).
 \label{eq:tpc42-scalar-operator-bound}
\end{equation}
Then, for every complex Hilbert space \(\Hh\) and every family
\((v_i)_{i\in I}\subset\Hh\),
\begin{equation}
 \sum_{j\in J}\left\|\sum_{i\in I}t_{ji}v_i\right\|_{\Hh}^2
 \le C^2\sum_{i\in I}\|v_i\|_{\Hh}^2.
 \label{eq:tpc42-hilbert-operator-bound}
\end{equation}
The constant is independent of \(\dim\Hh\), including when \(\Hh\) is
infinite dimensional.
\end{lemma}

\begin{proof}
All the vectors involved lie in a separable closed subspace.  Expand them
against an orthonormal basis \((e_\nu)\), apply
\eqref{eq:tpc42-scalar-operator-bound} to the scalar coefficients
\(\langle v_i,e_\nu\rangle\), and sum in \(\nu\).  Parseval and Tonelli give
\eqref{eq:tpc42-hilbert-operator-bound} with the same constant.
\end{proof}

For example, tensorizing the classical multiplicative large sieve
\citep{MontgomeryVaughan1973} gives, for finitely supported
\(v_n\in\Hh\),
\begin{equation}
 \sum_{r\le R}\frac r{\varphi(r)}
 \sum_{\chi\bmod r}^{*}
 \left\|\sum_{n\le Z}v_n\chi(n)\right\|_{\Hh}^{2}
 \ll (Z+R^2)\sum_{n\le Z}\|v_n\|_{\Hh}^{2}.
 \label{eq:tpc42-hilbert-multiplicative-large-sieve}
\end{equation}
The usual dyadic decomposition of initial intervals, now using
Cauchy--Schwarz in \(\Hh\), also gives the maximal benchmark
\begin{equation}
 \sum_{r\le R}\frac r{\varphi(r)}
 \sum_{\chi\bmod r}^{*}
 \max_{z\le Z}
 \left\|\sum_{n\le z}v_n\chi(n)\right\|_{\Hh}^{2}
 \ll (Z+R^2)(\log(2Z))^2
 \sum_{n\le Z}\|v_n\|_{\Hh}^{2}.
 \label{eq:tpc42-hilbert-maximal-large-sieve}
\end{equation}
These are standard tensor consequences of scalar large--sieve estimates;
they are recorded as a benchmark, not claimed as a new vector--valued
large sieve.

\subsection{Simultaneous Bessel transfer of the TPC--41 vector}

Let \(\mathfrak M_X\) be the TPC--41 family obtained by taking one prime
from each of three fixed, pairwise disjoint, fixed--proportion intervals
at scale \(X^{133/400}\).  For \(M\in\mathfrak M_X\), write
\begin{equation}
 \mathcal U_M=(\Z/M\Z)^{\times},
 \qquad K_M=X/M,
 \label{eq:tpc42-reduced-residue-space}
\end{equation}
and, for a family of bounded--variation weights \(W=W_X\), supported in
a fixed compact subinterval of \((0,\infty)\) and with uniformly bounded
BV norm, define
\begin{equation}
 S_{M,W}(a)
 =\sum_{\substack{n\equiv a\pmod M}}
   \mu(n)W(n/X),
 \qquad a\in\mathcal U_M.
 \label{eq:tpc42-scalar-residue-vector}
\end{equation}
TPC--41 proves that, for every fixed \(0<\theta<1/200\) and every fixed
finite collection of such uniformly bounded weights, all but
\(o(|\mathfrak M_X|)\) moduli satisfy simultaneously
\begin{equation}
 \|S_{M,W}\|_{\ell^2(\mathcal U_M)}^2
 \ll_{\theta,W}\frac{XK_M}{(\log X)^\theta}.
 \label{eq:tpc42-inherited-scalar-residue-bound}
\end{equation}
See \citet{WangTPC41}; the underlying centered progression input is due
to \citet{KlurmanMangerelTeravainen2023}.

\begin{proposition}[Simultaneous Bessel postcomposition]
\label{prop:tpc42-simultaneous-bessel-transfer}
Fix a modulus \(M\) for which
\eqref{eq:tpc42-inherited-scalar-residue-bound} holds.  Let \(\Hh_M\) be
any complex Hilbert space and let
\begin{equation}
 B_M:\ell^2(\mathcal U_M)\longrightarrow\Hh_M
 \label{eq:tpc42-bessel-operator}
\end{equation}
be any bounded linear map.  Then
\begin{equation}
 \boxed{
 \|B_MS_{M,W}\|_{\Hh_M}^2
 \ll_{\theta,W}
 \|B_M\|_{\op}^{\,2}
 \frac{XK_M}{(\log X)^\theta}.}
 \label{eq:tpc42-bessel-postcomposition-bound}
\end{equation}
The same good modulus works simultaneously for every choice of
\(\Hh_M\) and \(B_M\); in particular, \(B_M\) may be selected after the
scalar good set has been fixed.

Equivalently, suppose that \(v_{M,a}\in\Hh_M\) has Bessel synthesis bound
\begin{equation}
 \left\|\sum_{a\in\mathcal U_M}z_av_{M,a}\right\|_{\Hh_M}^2
 \le \beta_M\sum_{a\in\mathcal U_M}|z_a|^2.
 \label{eq:tpc42-residue-bessel-family}
\end{equation}
Then
\begin{equation}
 \left\|
  \sum_{\substack{n\ge1\\(n,M)=1}}
  \mu(n)W(n/X)v_{M,n\bmod M}
 \right\|_{\Hh_M}^2
 \ll_{\theta,W}
 \beta_M\frac{XK_M}{(\log X)^\theta}.
 \label{eq:tpc42-residue-synchronous-output}
\end{equation}
\end{proposition}

\begin{proof}
The first assertion is simply
\(\|B_MS_{M,W}\|\le\|B_M\|_{\op}\|S_{M,W}\|_2\), followed by
\eqref{eq:tpc42-inherited-scalar-residue-bound}.  No exceptional set is
introduced by this deterministic inequality.  For the second assertion,
take \(B_Me_a=v_{M,a}\).  Then
\eqref{eq:tpc42-residue-bessel-family} says
\(\|B_M\|_{\op}^2\le\beta_M\), and expanding \(B_MS_{M,W}\) gives
\eqref{eq:tpc42-residue-synchronous-output}.
\end{proof}

For the terminal logarithmic atom, take a fixed smooth compactly
supported \(\Phi\) and
\[
 W_X(t)=\frac{\log(Xt)}{\log X}\Phi(t).
\]
This family has uniformly bounded variation.  Multiplying
\eqref{eq:tpc42-bessel-postcomposition-bound} by \((\log X)^2\) yields
\begin{equation}
 \left\|
 B_M\left(
  \sum_{\substack{n\equiv a\pmod M}}
  \mu(n)\log n\,\Phi(n/X)
 \right)_{a\in\mathcal U_M}
 \right\|_{\Hh_M}^{2}
 \ll_{\theta,\Phi}
 \|B_M\|_{\op}^{\,2}
 \frac{XK_M(\log X)^2}{(\log X)^\theta}.
 \label{eq:tpc42-terminal-bessel-postcomposition}
\end{equation}

The qualifier \emph{simultaneous} is useful here.  Even if the physical
projective decomposition contains many residue masks or Hilbert output
layers, no union bound over those layers is needed when they are all
bounded postcompositions of the same vector \(S_{M,W}\).  Their entire
cost is the appropriate operator or Bessel norm.

\subsection{The precise stopping point}

\Cref{prop:tpc42-simultaneous-bessel-transfer} applies when the
coefficient attached to \(n\) is residue synchronous:
\begin{equation}
 G_M(n)=W(n/X)v_{M,n\bmod M}.
 \label{eq:tpc42-residue-synchronous-field}
\end{equation}
It does not estimate
\begin{equation}
 \sum_n\mu(n)G_M(n)
 \label{eq:tpc42-preweighted-vector-sum}
\end{equation}
for an arbitrary field \(G_M(n)\in\Hh_M\) that changes direction among
integers in the same residue fiber.  Inserting \(G_M(n)\) before applying
the scalar theorem changes its coefficient sequence, whereas
postcomposition changes only its already formed residue output.

This also explains why the terminology \emph{vector--valued KMT theorem}
would be inaccurate.  The logarithmic saving in
\eqref{eq:tpc42-inherited-scalar-residue-bound} first uses the scalar
multiplicativity and pretentious structure in KMT to identify and
subtract one most--correlated character; TPC--41 then removes that
rank--one term on average over the declared modulus family.  A general
map \(n\mapsto G_M(n)\) has no such multiplicative structure.  Only the
postcomposition step is a dimension--free tensorization.  The
within--fiber directional component left after the canonical residue
projection therefore requires a genuinely different arithmetic input at
the desired saving; it is not an omitted application of
\eqref{eq:tpc42-hilbert-multiplicative-large-sieve}; compare the exact
factorization criterion in \cref{thm:tpc42-factorization-criterion}.
